\documentclass[../../main.tex]{subfiles}
\graphicspath{{../../image/}{../../figure/}}
\begin{document}

\section{注记}\label{section:ch10-注记}

本章广义特征线理论的叙述遵循 Dafermos [19]。微分包含一般理论的系统阐述可见于 Filippov 的专著 [1]。早期将广义特征线（针对标量守恒律）作为古典特征方程在 Filippov 意义下的解引入的论文是 Wu [1]。另见 Hörmander [1]。Glimm 与 Lax [1] 采用了另一种广义特征线定义，即要么以古典特征速度、要么以激波速度传播的 Lipschitz 曲线，作为一族“近似特征线”的极限而构造。由\cref{theorem:ch10-广义特征线速度约束}可知，这两种定义密切相关。

分界（divide）的概念由 Dafermos [21] 引入。

尽管此处广义特征线是在 $BV$ 解的框架下讨论的，但它们完全可以在更广泛的 $L^{\infty}$ 弱解框架中定义和使用：在该框架中，\cref{definition:ch10-广义i特征线}仍有意义。在应用中遇到的一类弱解——它们未必属于 $BV$，但通过广义特征线进行处理仍然有效——包括其特征速度属于 $BV$ 的解，以及对一切 $x\in(-\infty,\infty)$、几乎所有 $t\in(0,\infty)$ 一侧极限 $U(x\pm,t)$ 均存在的解。

多维情形下的广义特征线由 Poupaud 与 Rascle [1] 讨论，其背景是具有间断系数的线性输运方程。

\end{document}